\chapter{Research and Existing Solutions}
\label{chap:research}
%=====================================================================

\section{Research Method / Approach}
The team investigated the domain in three strands:
\begin{itemize}
\item \textbf{Domain and need analysis:} epidemiology of cognitive decline,
ageing in place, and caregiver burden
\cite{who2024,statista2022,alz2023,mogle2023,marchant2011,kasper2015,harris2006,arno1999}.
\item \textbf{Technical literature:} recent work on assistive object retrieval
for older adults, long egocentric video question answering, and spatio-temporal
scene memory for agents
\cite{daaam2026,mempal2025,r4dsg2026,spatialmem2026,supermemory2026}.
\item \textbf{Product comparison:} features of existing assistive-recall tools
against the project's requirements.
\end{itemize}
Sources were selected by relevance to (a) object-location recall for people
with cognitive impairment, (b) persistent scene or event memory from
first-person streams, and (c) natural-language querying of that memory. Recent
venues (IUI 2025, CVPR 2026, ACM MM 2026) were prioritised for the technical
strand. In Iteration~2 the team added a fourth strand: its own component
experiments. Where a choice between methods could not be settled from the
literature, it was measured on the project's own data
(Section~\ref{sec:experiments}).

\section{Existing Solutions / Products}
\textbf{MemPal \cite{mempal2025}} (Maniar et al., IUI 2025) is a wearable,
voice-based memory assistant developed at the MIT Media Lab for older adults.
A multimodal large language model turns visual context from a wearable camera
into an automatically generated, real-time text diary of the wearer's
activities; the user then asks spoken questions to find misplaced objects or to
recall whether an action was completed. The authors tested it with 15 older
adults in their own homes and report better object-finding performance with
audio assistance than with no aid. The project page also describes a companion
app that tracks completed actions for remote monitoring by caregivers and
produces a summarised activity diary \cite{mempalpage}.

\textbf{R4DSG \cite{r4dsg2026}} (ACM MM 2026) targets object-centric questions
over long egocentric video, such as where an item was moved or when it last
changed state. It argues that caption-based memories rarely preserve persistent
object identity, and that earlier 3D scene-graph methods assume stronger
geometry (depth, poses, or reconstructions) than free-motion wearable RGB video
provides. Instead it converts video into compact, queryable memory entries
indexed by time, place, persistent objects, anchor-relative change, and local
interaction context, using RGB-only components for segmentation, temporal
propagation, and relative 3D lifting. It is positioned as a memory substrate
for wearable assistants; the descriptions we reviewed do not include a
user-facing assistive application or caregiver access.

\textbf{DAAAM \cite{daaam2026}} (Gorlo, Schmid and Carlone, CVPR 2026) is a
real-time spatio-temporal memory for embodied agents. Given RGB-D video with
camera poses, it segments and tracks scene fragments, captions selected frames
and segments with a localised captioning model using batch processing for
speed, and organises the results into a hierarchical 4D scene graph that a
tool-calling agent queries. It is evaluated on spatio-temporal question
answering (NaVQA and a longer-duration extension, OC-NaVQA). It targets robots
and embodied agents, not people with cognitive impairment.

\textbf{CrossSense} is a London-based smart-glasses system for people with
early-stage dementia. Its AI companion, Wispy, observes the wearer's
surroundings and gives proactive audio and visual prompts that guide everyday
tasks such as making tea or getting dressed
\cite{crosssense-techradar,crosssense-bigissue}. It won the Longitude Prize on
Dementia, and a University of Sussex team tested it with 23 pairs of people
living with dementia and their carers \cite{crosssense-bigissue}. The public
descriptions we found emphasise real-time guidance; they do not describe a
queryable long-term memory or caregiver query access. 

\textbf{Room-level mapping methods.} Aegis borrows the online trajectory clustering and revisit-merging idea from Hydra's spatial mapping work \cite{hydra2022}; it uses an HOV-SG-inspired separation of room-level semantic typing from lower-level object memory \cite{hovsg2024}. These are design inspirations adapted to sparse RGB-derived trajectories and detected landmarks, not reproductions of either complete method.

\textbf{Related egocentric-memory research.} Beyond direct competitors, the
wider literature motivates the design. Ego4D \cite{grauman2022} established
large-scale egocentric benchmarks, including episodic-memory tasks. EgoLife
\cite{egolife2025} extends the problem to an egocentric life assistant with
long-context tasks such as event recall and habit understanding, and AMEGO
\cite{amego2024} builds an active memory of recurring locations and object
interactions so that a very long video need not be reprocessed in full (as
summarised in the related-work section of \cite{r4dsg2026}). SpatialMem
\cite{spatialmem2026} organises memory hierarchically from egocentric RGB video
as anchors, objects, and text, and answers queries by walking that structure.
SuperMemory-VQA \cite{supermemory2026} is a benchmark of 4,853 question--answer
pairs on long-horizon egocentric memory, including object and location memory.
The team now uses EgoLife (participant A1\_JAKE) for its multi-day persistence
study (Section~\ref{sec:egolife}); Ego4D was assessed and rejected for FYP-2
(NO-GO, 0/100 readiness).

\section{Comparison of Existing Solutions}
Table~\ref{tab:comparison} compares the four systems with Aegis on
project-relevant criteria. ``Not described'' means the sources reviewed do not
state it, not that the system lacks it.


\begin{table}[H]
\centering\scriptsize
\setlength{\tabcolsep}{3pt}
\caption{Comparison of existing solutions with Aegis.}
\label{tab:comparison}
\begin{tabularx}{\textwidth}{@{}>{\bfseries}L{1.9cm} Y Y Y Y Y@{}}
\toprule
\rowcolor{hdrfill}\textbf{Criterion} & \textbf{MemPal \cite{mempal2025}} &
\textbf{CrossSense} & \textbf{R4DSG \cite{r4dsg2026}} &
\textbf{DAAAM \cite{daaam2026}} & \textbf{Aegis (FYP-1)}\\
\midrule
Target users & Older adults finding objects at home & Early-stage dementia &
Research: wearable assistants & Research: embodied agents & Broad cognitive
impairment, plus caregiver\\
Sensing input & Wearable camera & Smart glasses & RGB-only wearable video &
RGB-D with poses & Phone RGB camera (worn), MJPEG\\
Memory form & Text activity diary & Not described & Anchor-relative 4D
scene-graph memory entries & Hierarchical 4D scene graph & Rooms $\rightarrow$ anchors $\rightarrow$ objects; state runs and events; per-home SQLite store\\
Temporal abstraction & Not described & Not described & Timelines of object
state changes & Timeline of object placements & Implemented: three-day detail, verified day summaries, week/month roll-up\\
Interface & Voice assistant & Proactive audio and visual guidance & Text video
question answering & Tool-calling agent & Patient and caregiver apps; text, on-device speech, object tiles\\
Caregiver support & Remote monitoring, summarised diary & Not described &
Not described & Not described & Caregiver app with history, rooms, enrolment; consent-scoped access planned\\
Scaling approach & Not described & Not described & Compact memory entries &
Batch captioning of selected frames & State runs (size independent of video length) plus compaction\\
\bottomrule
\end{tabularx}
\end{table}

\section{Limitations and Gaps}
From the comparison, five gaps motivate Aegis:
\begin{enumerate}
\item \textbf{Spatial grounding for assistive use.} MemPal's memory is a text
activity diary \cite{mempal2025}; anchor-relative or metric object locations
are not described. DAAAM's 4D scene graph requires RGB-D input and camera poses
\cite{daaam2026}, which a wearable RGB camera may not supply. R4DSG shows that
anchor-relative memory can be built from RGB-only wearable video
\cite{r4dsg2026} but is not packaged as an assistive product.
\item \textbf{No progressive temporal abstraction in assistive systems.}
CrossSense is described as real-time guidance, and MemPal as a diary. R4DSG and
DAAAM keep timelines of object changes and placements, but a tiered policy that
condenses old history into progressively coarser summaries is not presented as
a design goal in any reviewed system. The room layer provides a natural coarse
spatial tier that may preserve location context through future temporal
compaction; retention of room identity through compaction is now verified by
construction (rooms are never compacted), but its effect on query accuracy is
unmeasured.
\item \textbf{Research systems are not assistive systems.} R4DSG and DAAAM show
that object-centric, anchor-relative, and hierarchical scene memory work for
question answering and for embodied agents, but they are not designed for
people with cognitive impairment, nor with speech access, evidence display, or
caregiver consent.
\item \textbf{Caregiver access is passive or absent.} MemPal offers remote
monitoring and a summarised diary; none of the reviewed systems describes
delegated, consent-scoped querying of the full history by a caregiver.
\item \textbf{No unified pipeline.} None of the reviewed systems combines
continuous egocentric ingestion, tiered temporal decay, and multimodal
natural-language grounding in one engineered pipeline.
\end{enumerate}
The gap, in one sentence: \emph{among the systems reviewed, none combines
persistent, spatially grounded egocentric memory with progressive temporal
abstraction and multimodal, evidence-based querying for both a person with
cognitive impairment and their caregiver.} Figure~\ref{fig:gap}
(Section~\ref{sec:proposed}) positions Aegis in this unoccupied space.

\section{Proposed Solution}
\label{sec:proposed}
Aegis is a five-stage pipeline with two access paths (Figure~\ref{fig:pipeline}).
The write path is continuous and the read path is on demand; both touch the
shared memory store.

\aegisfigtall{ch2-pipeline}
  {The Aegis pipeline at the end of FYP-1: continuous write path, on-demand
   read path, and the shared memory store. Dashed: planned.}
  {fig:pipeline}

\textbf{Stage walkthrough.} Capture and ingestion (FR1, FR2) turns first-person
video into a frame stream and selects keyframes to limit cost. The
spatio-temporal memory engine (FR3) detects and tracks objects, anchors them
relative to stable parts of the scene, and groups observations into events.
Multi-resolution abstraction (FR4) keeps recent history detailed and folds
older history into compact summaries; compaction runs in the background. At
query time the interface (FR6, FR7) accepts speech or text, the reasoning agent
(FR5) retrieves the latest states from the memory, and a template composes
the answer from the retrieved state; a language model is optional and off by
default (temporal and activity questions are declined until Iteration~3). Patient and caregiver access (FR7, FR8) present the same
memory through two permissioned views. The write and read paths are drawn
separately because ingestion is continuous while queries are on demand;
queries do not flow \emph{through} the compression stage but read the store it
maintains. The design decisions behind the stages are documented in
Section~\ref{sec:decisions}.

Figure~\ref{fig:gap} places Aegis against the reviewed systems on the two axes
that define the gap: how finely space is represented, and how far memory
extends and is abstracted over time.

\aegisfig[0.88\textwidth]{fig-2.2}
  {Gap positioning: spatial grounding versus temporal memory. Placements are
   the team's qualitative judgement from published descriptions
   (Table~\ref{tab:comparison}).}
  {fig:gap}

\section{Complex Computing Problem}
Aegis addresses a complex computing problem. The characteristics below follow
the attribute framework used in the project's proposal submission. The concrete
constraint is deliberate: retain structured, spatially grounded memory for three
days, then convert older spatial observations into text summaries. This creates a
fidelity-versus-storage trade-off rather than claiming indefinite spatial
persistence.

\begin{table}[H]
\centering\footnotesize
\caption{Complex computing problem attributes.}
\label{tab:complex}
\begin{tabularx}{\textwidth}{@{}c L{3.0cm} Y Y@{}}
\toprule
\rowcolor{hdrfill}\textbf{\#} & \textbf{Attribute} & \textbf{How Aegis exhibits it} & \textbf{Evidence (FYP-1)}\\
\midrule
1 & Conflicting requirements & Preserving fine spatial fidelity over time
conflicts with memory scalability. Maintaining a comprehensive event history
conflicts with low-latency, low-cognitive-load responses. & Measured trade-off: snapshot interval 1 keeps exact geometry at 1.5--1.9$\times$ compression; interval 0 gives 5.7--7.2$\times$ with centre p95 error 0.22--0.36 scene units (Section~\ref{sec:compaction-audit}).\\
2 & No obvious solution; requires conceptual modelling & No unified pipeline
integrates continuous egocentric ingestion, tiered temporal decay, and
multimodal grounding, so an original conceptual memory model is required. & --\\
3 & In-depth knowledge based on first principles & Requires geometric vision,
probabilistic state estimation (tracking objects under uncertainty), and
semantic retrieval. & --\\
4 & Infrequently encountered issues & Continuous, unconstrained first-person
streams with dynamic state changes, rather than static datasets. & --\\
5 & Outside standard practice & No engineering standard governs continuous
episodic memory abstraction. & --\\
6 & Diverse stakeholders & A primary user needing non-intrusive assistance and
a caregiver needing verifiable oversight. & --\\
7 & Significant consequences & Retrieval failures affect daily independence and
caregiver intervention burden. & --\\
8 & Interdependent sub-problems & Changes in tracking or compression
granularity propagate into query accuracy. & Changing the detector (YOLOE $\rightarrow$ SAM~3) changed which objects exist in memory at all (0/36 $\rightarrow$ 34/42 shoes); the room-linking radius changed cross-session identity (Section~\ref{sec:egolife}).\\
9 & Functional requirement categories identifiable & Requirement categories are
drawn from established assistive cognition literature even though the
architecture is novel. & --\\
\bottomrule
\end{tabularx}
\end{table}

\textbf{Why this is not routine engineering.} A concrete example: if compaction
discards the exact shelf an object last rested on, a later query for that
object fails even though capture and tracking worked. Compaction policy,
tracking quality, and retrieval design therefore cannot be tuned independently.
\section{Technical Challenges}
\begin{ltab}\scriptsize
\begin{longtable}{@{}>{\bfseries}L{0.9cm}L{2.3cm}L{3.7cm}L{3.7cm}L{3.7cm}@{}}
\caption{Technical challenges, planned mitigations and what FYP-1 found.}\label{tab:challenges}\\
\toprule
\rowcolor{hdrfill}\textbf{ID} & \textbf{Challenge} & \textbf{Why it is hard} &
\textbf{Planned mitigation} & \textbf{What we found in FYP-1}\\
\midrule
\endfirsthead
\toprule
\rowcolor{hdrfill}\textbf{ID} & \textbf{Challenge} & \textbf{Why it is hard} &
\textbf{Planned mitigation} & \textbf{What we found in FYP-1}\\
\midrule
\endhead
\bottomrule
\endfoot
TC1 & Persistent object identity in cluttered, dynamic homes & Objects are
occluded, moved, and look similar; identity may break or duplicate. &
Anchor-relative tracking; re-identification across sightings; confidence scores
on identity. & Re-ID moved from HSV histograms to DINOv3 embeddings; multi-instance identity and occlusion permanence are tested (\texttt{test\_association}, \texttt{test\_scene\_graph}); not yet measured on ground truth.\\
TC2 & Abstraction trade-off & Compression may discard the fine detail the user
needs. & Tiered retention; importance-aware summarisation; test accuracy
before and after compaction. & Compaction verifies each summary against its source rows before deleting detail; a failed day keeps all rows.\\
TC3 & Egocentric video quality & Head motion blur, narrow view, lighting
changes. & Keyframe selection that scores sharpness and novelty. & Criterion keyframes reject dark ($<$20), bright ($>$245), blurred (Laplacian $<$15) and duplicate ($<$0.8\% change within 2\,s) frames.\\
TC4 & Efficient VLM use & VLM calls are costly and slow at video rates. &
Process keyframes only; cache descriptions; batch captioning, as in
\cite{daaam2026}. & VLM removed from the answer path; GPU used only for reconstruction and detection; Modal warm-down 2\,s.\\
TC5 & Query grounding and hallucination & Free-form questions may be vague,
relative (``before lunch''), or unanswerable. & Retrieval-first answering;
require evidence; explicit abstention. & Closed vocabulary and templates; the engine declines unknown items, unobserved items, temporal and activity questions.\\
TC6 & Low-load interaction for impaired users & Long or complex answers can
confuse. & Short answers with a single piece of visual evidence. & Answers are one sentence with up to three photos and a labelled certainty badge.\\
TC7 & Privacy and consent & Continuous capture of private spaces and
bystanders. & Local-first storage, minimal raw retention, access control. & Not yet addressed: no authentication or encryption (P1.7, P1.8).\\
TC8 & Limited validation & Small-scale tests may not generalise to real
dementia patients. & Report limitations openly; test with scripted scenarios
and healthy volunteers. & Unchanged; honest limits reported for every measurement.\\
TC9 & Dataset not yet selected & No reviewed dataset is guaranteed to match the indoor, caregiver-oriented setting. & Select a suitable dataset or subset in FYP-2; a relevant portion of Ego4D is the current candidate, with annotation/adaptation documented before evaluation. & EgoLife chosen for the multi-day study; Ego4D NO-GO.\\ 
TC10 & Spatial reasoning from RGB-only wearable video & No depth sensor or
external pose; the camera moves freely. & Use raw VGGT-$\Omega$ coordinates
and floor-plane geometry without claiming calibrated metric accuracy. & A gravity frame from a RANSAC floor fit \cite{ransac1981} gives ``above / under / on floor'' without metric scale.\\
TC11 & Room segmentation and cross-session identity & Trajectory clustering may
split/merge regions after revisits; furniture moves can change landmark evidence.
& Evaluate room assignments and transitions, test registry matching/blending,
and report failures. & Room signatures normalised (\texttt{61dda40}); identity still unstable across sessions (TC12).\\
TC12 & Cross-session room identity under scale-free reconstruction & \texttt{zone\_link\_radius} (1.5) is an absolute distance, but VGGT scale differs per reconstruction. & Scale-relative linking using the gravity frame's camera-height unit (Iteration~3). & On two EgoLife segments one produced four rooms and the next one room, matched to the wrong home room.\\
\end{longtable}
\end{ltab}

\section{Research Findings and Design Decisions}
\label{sec:decisions}
\begin{ltab}\footnotesize
\begin{longtable}{@{}L{7.2cm}L{7.2cm}@{}}
\caption{Research findings and the design decisions they drive.}\label{tab:decisions}\\
\toprule
\rowcolor{hdrfill}\textbf{Research finding (source)} & \textbf{Design decision in Aegis}\\
\midrule
\endfirsthead
\toprule
\rowcolor{hdrfill}\textbf{Research finding (source)} & \textbf{Design decision in Aegis}\\
\midrule
\endhead
\bottomrule
\endfoot
MemPal builds a text diary from a wearable camera and supports voice object
retrieval and caregiver monitoring \cite{mempal2025,mempalpage}. & Natural-language,
accessible querying and an activity-event memory as the interaction baseline,
extended with anchor-relative spatial detail and delegated caregiver querying.\\
R4DSG builds compact memory entries indexed by time, place, persistent objects,
and anchor-relative change from RGB-only wearable video \cite{r4dsg2026}. &
Room--anchor--object scene graph with timestamped events; raw VGGT coordinates
are used without a manual scale or axis transform (Section~\ref{sec:constraints}).\\
DAAAM needs RGB-D and camera poses for its metric 4D scene graph
\cite{daaam2026}. & Aegis uses raw VGGT-$\Omega$ coordinates as its spatial
reference and floor-plane room geometry, while treating scale and room boundaries
as approximate rather than calibrated metric measurements.\\
DAAAM organises memory as a hierarchical scene graph with a timeline of
placements \cite{daaam2026}. & Hierarchical organisation of environment, areas,
anchors, and objects, with a timeline of placements.\\
DAAAM selects frames and segments for captioning and uses batch processing to
make inference fast enough for online use \cite{daaam2026}. & Keyframe
optimisation in ingestion (Milestone M1).\\
A tool-calling agent can query a spatio-temporal memory for answers
\cite{daaam2026}. & Planned at mid as a voice agent; replaced in Iteration~2 by a
closed-vocabulary engine for grounding (deviation D1).\\
R4DSG stores compact entries instead of raw graph sequences
\cite{r4dsg2026}. & Progressive high-level summaries as the older tier of
memory.\\
Dementia places a disproportionate burden on unpaid caregivers
\cite{kasper2015}, and informal care is costly \cite{arno1999,chari2015}. &
Delegated caregiver application with history and evidence; consent-based access planned.\\
Some people with early-stage Alzheimer's disease live alone and want to cope
independently \cite{harris2006}. & Accessible interaction with simple patient and caregiver views.\\
Hydra performs online spatial clustering and revisit merging \cite{hydra2022}; HOV-SG organizes open-vocabulary scene graphs hierarchically \cite{hovsg2024}. & Estimate room regions from wearer trajectory, merge revisits, type rooms from landmark votes, and keep architectural landmarks out of the object timeline.\\
\rowcolor{rowfill}\multicolumn{2}{@{}l}{\textit{Decisions added in Iteration 2, from the team's own experiments}}\\
YOLOE found 0/36 shoes and 0/19 hands; SAM~3 \cite{sam3} found 34/42 and 19/19 (team benchmark, Section~\ref{sec:experiments}). & SAM~3 is the default detector for recorded video; YOLOE \cite{yoloe} stays available for live use.\\
Appearance must survive occlusion and viewpoint change between sightings. & DINOv3 \cite{dinov3} embeddings for re-identification instead of colour histograms.\\
Per-frame timelines grow with frames $\times$ objects (team memory audit). & Store debounced \textbf{state runs} instead; graph size independent of video length.\\
A free-form language model can produce a plausible but unsupported location (TC5). & Closed vocabulary (525 items) and template answers; an LLM only as an optional parser (P1.9).\\
VGGT-$\Omega$ places the origin at the first camera; ``up'' is unknown. & Estimate gravity from a RANSAC floor fit \cite{ransac1981}; store it beside the raw coordinates.\\
JSON graphs cannot merge sessions or survive restarts efficiently. & Per-home SQLite store with stable room, anchor and object identities.\\
Class-level queries cannot distinguish two mugs (EgoLife tiers). & Ranked top-3 candidates with photos and an explicit ambiguity flag; personal-item enrolment.\\
\end{longtable}
\end{ltab}


\section{Model and Method Selection Experiments}
\label{sec:experiments}
\textbf{Object detection.} Open-vocabulary detection decides which objects can
ever enter memory. The team hand-labelled a small benchmark (10 frames,
960$\times$540) from five different videos --- a shoe cabinet, a kitchen, a
bathroom/waste area, laundry and a dish rack --- with 116 objects, drawn by
inspecting raw frames rather than pipeline output. Detections were scored
class-aware at IoU $\geq$ 0.30 after per-class non-maximum suppression (IoU
0.7). Table~\ref{tab:detector-ab} and Figure~\ref{fig:detector-ab} compare five
variants on a Modal L4 GPU.

\begin{table}[H]
\centering\small
\caption{Detector comparison on the ten-frame benchmark (\texttt{docs/detection\_benchmark.md}).}
\label{tab:detector-ab}
\begin{tabularx}{\textwidth}{@{}Y r r r r r r r@{}}
\toprule
\rowcolor{hdrfill}\textbf{Variant} & \textbf{Prompts} & \textbf{P} & \textbf{R} & \textbf{F1} & \textbf{FP/img} & \textbf{Shoes} & \textbf{Hands}\\
\midrule
YOLOE-11l generic (old default) & 93 & 0.357 & 0.056 & 0.097 & 0.9 & 0/36 & 0/19\\
YOLOE-11l compact & 18 & 0.556 & 0.056 & 0.102 & 0.4 & 0/36 & 0/19\\
YOLO-World m compact & 18 & 0.460 & 0.191 & 0.270 & 2.0 & 11/36 & 0/19\\
SAM~3 compact, threshold 0.5 & 18 & 0.326 & 0.798 & 0.463 & 14.7 & 27/36 & 19/19\\
\textbf{SAM~3 compact, class thresholds, NMS 0.7 (shipped)} & 18 & \textbf{0.718} & \textbf{0.724} & \textbf{0.721} & 3.3 & 34/42 & \textbf{19/19}\\
\bottomrule
\end{tabularx}
\end{table}

\aegisfig{ch2-detector-ab}
  {Precision, recall and F1 of five detector variants; false positives per
   image above each group. 10 hand-labelled frames, 960$\times$540,
   five environments, 116 objects; class-aware IoU $\geq$ 0.30; Modal L4.}
  {fig:detector-ab}

YOLOE, the previous default, found none of the 36 shoes or 19 hands.
YOLO-World found 11 shoes but no hands. SAM~3 \cite{sam3} with an 18-class
prompt set and category-level thresholds (shoes 0.5, hand 0.4, cooking pot 0.6,
trash bin 0.4, food container 0.9, others 0.8) reached F1 0.721, and was
adopted. The thresholds are category priors, not tuned per video. Every
alternative tried (prompt-word variants, NMS sweeps, cross-class suppression,
flip test-time augmentation) was rejected when it regressed any video or gained
less than 0.005. The benchmark is small, so the result guides the choice of
detector; it is not a general accuracy claim.

\textbf{Memory representation and long-term persistence.} Two further
experiments informed the memory design and are reported in full in
Chapter~\ref{chap:testing}: the memory-compaction audit
(Section~\ref{sec:compaction-audit}), which motivated state runs, and the
seven-day EgoLife study (Section~\ref{sec:egolife}), which exposed technical
challenge TC12.
